\section{Phân tích và so sánh các giải pháp tương tự}
Trong những năm gần đây, nhiều nền tảng trực tuyến đã được phát triển nhằm hỗ trợ người dùng
trong việc tìm kiếm địa điểm ăn uống, đánh giá chất lượng quán ăn và định hướng lựa chọn phù
hợp với nhu cầu cá nhân. Tuy nhiên, các giải pháp hiện tại chủ yếu tập trung vào từng khía cạnh
riêng lẻ của bài toán ăn uống và chưa cung cấp một giải pháp tích hợp toàn diện cho việc xây
dựng và tối ưu hóa các hành trình ăn uống (food tour). Phần này tiến hành phân tích một số giải
pháp tiêu biểu nhằm làm rõ ưu điểm, hạn chế và khoảng trống mà đề tài hướng đến giải quyết.
\subsection{Google Maps}
\subsubsection{Mô tả tổng quan}
Google Maps là một ứng dụng bản đồ phổ biến, cung cấp các chức năng tìm kiếm địa điểm, dẫn
đường và hiển thị thông tin cơ bản về các quán ăn như vị trí, giờ mở cửa và đánh giá từ người
dùng. Ứng dụng hỗ trợ dẫn đường theo nhiều phương tiện di chuyển khác nhau và được sử dụng
rộng rãi trong đời sống hằng ngày.
\subsubsection{Ưu điểm}
\begin{itemize}[label=\textbullet]
    \item Cung cấp dữ liệu bản đồ và định vị có độ chính xác cao.
    \item Hỗ trợ tìm kiếm quán ăn theo vị trí địa lý và dẫn đường trực tiếp.
    \item Có lượng người dùng lớn, dữ liệu đánh giá phong phú và thường xuyên được cập nhật.
    \item Giao diện thân thiện, dễ sử dụng trên nhiều thiết bị.
\end{itemize}
\subsubsection{Hạn chế}
\begin{itemize}[label=\textbullet]
    \item Chỉ tập trung vào việc tìm kiếm và dẫn đường cho từng địa điểm riêng lẻ.
    \item Không hỗ trợ trực tiếp việc xây dựng và quản lý các food tour gồm nhiều quán ăn.
    \item Việc sắp xếp thứ tự ghé thăm nhiều địa điểm trong cùng một hành trình còn mang tính thủ
công và phụ thuộc vào người dùng.
    \item Không có cơ chế chia sẻ food tour như một thực thể độc lập để người dùng khác có thể tìm
kiếm và tái sử dụng.
\end{itemize}

\subsection{Foody}
\subsubsection{Mô tả tổng quan}
Foody là ứng dụng tập trung vào việc cung cấp thông tin, đánh giá và gợi ý quán ăn dựa
trên trải nghiệm của cộng đồng người dùng. Các ứng dụng này cho phép người dùng tìm kiếm
quán ăn theo loại hình, khu vực hoặc mức giá, đồng thời xem hình ảnh và nhận xét chi tiết.
\subsubsection{Ưu điểm}
\begin{itemize}[label=\textbullet]
    \item Cung cấp thông tin chi tiết về quán ăn, bao gồm hình ảnh, thực đơn và đánh giá.
    \item Hỗ trợ gợi ý quán ăn dựa trên xu hướng hoặc đánh giá của cộng đồng.
    \item Giao diện trực quan, phù hợp với nhu cầu tham khảo nhanh.
    \item Dễ dàng tiếp cận đối với người dùng phổ thông.
\end{itemize}
\subsubsection{Hạn chế}
\begin{itemize}[label=\textbullet]
    \item Tập trung vào việc đánh giá và gợi ý quán ăn đơn lẻ, chưa hỗ trợ xây dựng food tour hoàn chỉnh.
    \item Không hỗ trợ tối ưu lộ trình di chuyển giữa nhiều quán ăn.
    \item Việc lưu trữ và quản lý các hành trình ăn uống cá nhân còn hạn chế.
    \item Chưa cung cấp cơ chế chia sẻ food tour một cách có hệ thống.
\end{itemize}

\subsection{TripAdvisor}

\subsubsection{Mô tả tổng quan}
TripAdvisor là một nền tảng trực tuyến quốc tế chuyên cung cấp thông tin, đánh giá và xếp hạng các địa điểm du lịch, bao gồm khách sạn, nhà hàng và các điểm tham quan.
Đối với lĩnh vực ẩm thực, TripAdvisor cho phép người dùng tìm kiếm quán ăn theo vị trí địa lý, loại hình ẩm thực và mức độ đánh giá, đồng thời tham khảo các nhận xét chi tiết từ cộng đồng người dùng trên toàn thế giới.

\subsubsection{Ưu điểm}
\begin{itemize}[label=\textbullet]
    \item Cung cấp lượng lớn dữ liệu đánh giá và nhận xét từ cộng đồng người dùng quốc tế.
    \item Hỗ trợ tìm kiếm và lọc quán ăn theo nhiều tiêu chí như loại ẩm thực, vị trí và mức độ phổ biến.
    \item Thông tin quán ăn được trình bày chi tiết, bao gồm hình ảnh, đánh giá và xếp hạng tổng thể.
    \item Phù hợp cho người dùng có nhu cầu tham khảo và khám phá ẩm thực tại các khu vực đô thị và du lịch.
\end{itemize}

\subsubsection{Hạn chế}
\begin{itemize}[label=\textbullet]
    \item Tập trung chủ yếu vào việc đánh giá và xếp hạng các quán ăn đơn lẻ, chưa hỗ trợ xây dựng food tour.
    \item Không cung cấp chức năng lập kế hoạch hành trình ăn uống gồm nhiều điểm dừng.
    \item Chưa hỗ trợ tối ưu lộ trình di chuyển giữa các quán ăn trong cùng một hành trình.
    \item Việc chia sẻ và tái sử dụng các hành trình ăn uống chưa được hỗ trợ như một thực thể độc lập.
\end{itemize}

\subsection{Tổng kết và nhận xét}
Qua phân tích các nền tảng tiêu biểu bao gồm nền tảng bản đồ (Google Maps), nền tảng đánh giá quán ăn nội địa (Foody) và nền tảng đánh giá du lịch quốc tế (TripAdvisor), có thể nhận thấy rằng các giải pháp hiện tại đều hướng đến việc hỗ trợ người dùng tìm kiếm và lựa chọn quán ăn. Tuy nhiên, mỗi nền tảng chỉ giải quyết một phần của bài toán tổng thể liên quan đến hành trình ăn uống.

Các nền tảng bản đồ có ưu thế về định vị và dẫn đường nhưng thiếu chức năng quản
lý hành trình ăn uống, trong khi các nền tảng đánh giá quán ăn cung cấp nhiều thông tin tham khảo nhưng chưa hỗ trợ tối ưu hóa lộ trình di chuyển.
\subsubsection{Đặc điểm chung của các nền tảng}
\begin{itemize}[label=\textbullet]
    \item Hỗ trợ tìm kiếm và hiển thị thông tin quán ăn.
    \item Cung cấp đánh giá hoặc phản hồi từ cộng đồng người dùng.
    \item Thiếu khả năng quản lý food tour như một thực thể độc lập.
\end{itemize}
\subsubsection{Những hạn chế còn tồn tại}
\begin{itemize}[label=\textbullet]
    \item Chưa có giải pháp tích hợp đầy đủ giữa tìm kiếm quán ăn, xây dựng food tour và tối ưu lộ trình.
    \item Việc lập kế hoạch cho các hành trình ăn uống nhiều điểm vẫn mang tính thủ công.
    \item Chưa khai thác hiệu quả khả năng chia sẻ và tái sử dụng food tour giữa các người dùng.
\end{itemize}
\subsubsection{Cơ hội phát triển và định hướng của đề tài}
Từ những hạn chế trên, có thể thấy còn tồn tại khoảng trống cho một hệ thống tích hợp hỗ trợ
người dùng xây dựng, quản lý và chia sẻ food tour một cách hiệu quả. Đề tài này hướng đến việc
giải quyết khoảng trống đó bằng cách phát triển một ứng dụng web tập trung vào food tour như
một đối tượng trung tâm, kết hợp với chức năng sắp xếp lộ trình di chuyển hợp lý và chia sẻ
công khai, từ đó nâng cao trải nghiệm người dùng trong việc khám phá ẩm thực đô thị.

